\phantomsection
\addcontentsline{toc}{section}{Abstract}
\section*{Abstract}

Modern municipal smart city grids monitor vital water distribution, electrical power, and air quality infrastructure using distributed Internet of Things (IoT) sensors. However, deploying edge sensing nodes in public spaces introduces severe cyber-physical and privacy vulnerabilities, including packet flooding bursts, physical value spoofing (e.g., impossible negative flow or out-of-scale AQI), and citizen surveillance through unencrypted geolocation and hardware identifiers. This project presents the design, specification, implementation, and empirical validation of \textbf{SIGS-RS (Secure IoT Gateway Simulation for Resource Sensing)}, an asynchronous edge gateway built with Python 3.12, FastAPI, and TimescaleDB that enforces a strict sequential architectural invariant: \emph{raw telemetry must traverse two-stage security middleware before reaching persistent storage}.

The gateway accepts telemetry via HTTP REST and secure MQTT over TLS 1.2+ (port 8883, $\le 512$~KB ceiling), buffering ingress into an in-memory \code{IngestionQueue} to decouple network sockets from relational writes. In Stage~1, the \textbf{Anomalous Packet Inspector} evaluates incoming packets against a sliding-window rate limit (50 packets per 10 seconds via in-memory timestamp deques or distributed Redis sorted sets) and physical validity equations; invalid packets are dropped immediately to an immutable audit log without reaching Stage~2. In Stage~2, the \textbf{Privacy Anonymizer} permanently purges geographic coordinates (\code{latitude}, \code{longitude}) and hardware identifiers (\code{mac_address}, \code{device_id}) under a fail-secure redaction policy. Sanitized records are committed to partitioned TimescaleDB hypertables protected by an adaptive persistence \textbf{Circuit Breaker} ($> 2,000$~ms trip threshold). System governance is enforced via SRS BR-3, which automatically locks out telemetry ingress across all gateway instances upon privacy guideline updates until formal audit signoff, while PostgreSQL triggers prevent \texttt{UPDATE} and \texttt{DELETE} tampering on \code{audit.security_log}.

The platform was verified through 580 automated tests (579 passed, 1 skipped due to the absence of an external broker). Empirical microbenchmarking across 10 to 1,000 sensors demonstrated sub-millisecond middleware processing ($p99 < 0.50$~ms), fully satisfying the SRS 15.0~ms requirement (REQ-PF-2). In accordance with software engineering audit standards, high-capacity distributed targets---specifically sustained 50,000 packets/second across 10,000 sensors (REQ-PF-1) and hardware-level WORM storage---are explicitly disclosed as unverified on the local host and categorized as production cluster prerequisites.

\vspace{1.5em}
\noindent\textbf{Keywords:} Secure IoT Gateway, Resource Sensing, Sliding-Window Rate Limiting, Privacy Anonymization, Circuit Breaker, Fail-Closed Security, Software Requirements Specification, TimescaleDB.
\clearpage
